\documentclass[11pt]{article}

\usepackage[a4paper,margin=1in]{geometry}
\usepackage{amsmath,amssymb,amsthm}
\usepackage{graphicx}
\usepackage[hidelinks]{hyperref}

\newtheorem{theorem}{Theorem}
\newtheorem{lemma}{Lemma}
\newtheorem{corollary}{Corollary}
\newtheorem{remark}{Remark}

\newcommand{\C}{\mathbb{C}}
\newcommand{\norm}[1]{\left\lVert #1 \right\rVert}
\newcommand{\hot}{\widehat{\otimes}_{\gamma}}
\emergencystretch=3em

\title{An Infinite-Dimensional Banach Algebra With A Nonzero Hyper-Commutator}
\author{Automatic math-discovery smoke test}
\date{June 29, 2026}

\begin{document}
\maketitle

\section*{Source Problem}

The source paper is Maysam Maysami Sadr, \emph{A Remark on Contractible
Banach Algebras of Operators}, arXiv:2211.03493 \cite{Sadr2022}.  The paper
defines a hyper-commutator for a Banach algebra \(A\) to be an element
\(M\in A\hot A\) such that
\[
  aM=Ma\qquad(a\in A),
\]
where \(A\) acts on \(A\hot A\) by multiplication on the left and right tensor
factors.  Immediately after the definition, the paper asks:
\begin{quote}
Does there exist an infinite-dimensional Banach algebra with a nonzero
hyper-commutator?
\end{quote}

\IfFileExists{figures/open_problem_crop.png}{
\begin{center}
\includegraphics[width=\linewidth]{figures/open_problem_crop.png}
\end{center}
}{}

\section*{Result Type}

This packet gives a \textbf{candidate full affirmative answer} to the question
as stated.  In fact the example can be chosen unital.

\section*{Proof Intuition}

The fastest way to make the two tensor-module actions agree is to put a large
square-zero ideal \(I\) inside a unital algebra.  If \(I^2=0\), then elements
of \(I\) cannot see the \(I\)-part of an algebra element when multiplied on
either side; only the scalar part of the unitization remains.  Therefore every
elementary tensor \(u\otimes v\) with \(u,v\in I\) is scaled by the same scalar
on the left and on the right.  Choosing \(u\) and \(v\) nonzero gives a
nonzero hyper-commutator.

\section*{The Example}

\begin{lemma}\label{lem:unitization}
Let \(E\) be a Banach space.  On the Banach space
\[
  A=\C\oplus_1 E,\qquad
  \norm{(\lambda,x)}=|\lambda|+\norm{x},
\]
define
\[
  (\lambda,x)(\mu,y)=(\lambda\mu,\lambda y+\mu x).
\]
Then \(A\) is a unital Banach algebra with identity \((1,0)\).  The closed
ideal \(I=\{0\}\oplus E\) satisfies \(I^2=0\).
\end{lemma}

\begin{proof}
This is the usual unitization of the Banach algebra \(E\) equipped with the
zero product.  Associativity follows directly from associativity of scalar
multiplication.  The element \((1,0)\) is a two-sided identity.  Finally,
\[
\begin{aligned}
  \norm{(\lambda,x)(\mu,y)}
  &=|\lambda\mu|+\norm{\lambda y+\mu x}  \\
  &\le |\lambda||\mu|+|\lambda|\norm{y}+|\mu|\norm{x} \\
  &\le (|\lambda|+\norm{x})(|\mu|+\norm{y}),
\end{aligned}
\]
so multiplication is jointly continuous and submultiplicative for the
\(\ell^1\)-sum norm.  If \(x,y\in E\), then
\[
  (0,x)(0,y)=(0,0),
\]
so \(I^2=0\).
\end{proof}

\begin{theorem}\label{thm:main}
There exists an infinite-dimensional unital Banach algebra with a nonzero
hyper-commutator.
\end{theorem}

\begin{proof}
Let \(E\) be any infinite-dimensional Banach space, for instance \(E=\ell^1\),
and form the Banach algebra \(A=\C\oplus_1 E\) from Lemma
\ref{lem:unitization}.  Choose nonzero vectors \(u,v\in E\), and set
\[
  p=(0,u),\qquad q=(0,v),\qquad M=p\otimes q\in A\hot A.
\]
The projective tensor norm is a cross norm, so
\[
  \norm{p\otimes q}_{\gamma}=\norm{p}\,\norm{q}>0.
\]
Hence \(M\ne 0\) in \(A\hot A\).

Let \(a=(\lambda,x)\in A\).  Since \(I^2=0\),
\[
  ap=(\lambda,x)(0,u)=(0,\lambda u)=\lambda p,
  \qquad
  qa=(0,v)(\lambda,x)=(0,\lambda v)=\lambda q.
\]
Using the standard outer bimodule actions on the projective tensor product,
\[
  aM=(ap)\otimes q=\lambda(p\otimes q)
  =p\otimes(\lambda q)=p\otimes(qa)=Ma.
\]
Thus \(M\) is a nonzero hyper-commutator for the infinite-dimensional Banach
algebra \(A\).
\end{proof}

\begin{corollary}
The question quoted above from arXiv:2211.03493 has an affirmative answer.
\end{corollary}

\section*{Scope And Limitations}

The construction is intentionally elementary and answers the literal question
as written.  It does not address stronger variants in which one also requires
the algebra to be semisimple, to be represented faithfully as an operator
algebra containing finite-rank operators, to have nonzero image under a
specific canonical operator-tensor map, or to be contractible.

\section*{Verification Notes}

The only points needing verification are:
\begin{enumerate}
\item \(A=\C\oplus_1 E\) is a Banach algebra.  This is checked in Lemma
\ref{lem:unitization}.
\item \(p\otimes q\ne 0\) in the completed projective tensor product.  This
follows from the cross-norm identity for elementary tensors, or equivalently
by applying bounded functionals detecting \(p\) and \(q\).
\item The hyper-commutator identity uses only \(I^2=0\), so it holds for every
\(a\in A\).
\end{enumerate}

The accompanying script
\texttt{code/check\_square\_zero\_hypercommutator.py} checks the displayed
left/right action identities in a finite-dimensional square-zero model.

\section*{Novelty And Literature Check}

Before promotion, the run registry, solution, attempt, and proof-gap indexes
were searched for arXiv:2211.03493, the source title, \emph{hyper-commutator},
\emph{nonzero hyper-commutator}, and related phrases.  No duplicate packet,
prior attempt, or proof-gap entry answering this question was found.

A bounded web search on June 29, 2026 for exact phrases including
\emph{nonzero hyper-commutator}, \emph{hyper-commutator Banach algebra}, and
the full question found the source arXiv page but no separate later answer.
The local parsed-source corpus likewise found the hyper-commutator question in
the source paper but no later explicit resolution.

\begin{thebibliography}{9}
\bibitem{Sadr2022}
Maysam Maysami Sadr,
\emph{A Remark on Contractible Banach Algebras of Operators},
arXiv:2211.03493, 2022.
\end{thebibliography}

\end{document}
